\documentclass[11pt]{article}
\usepackage[margin=1in]{geometry}
\usepackage{enumitem}
\usepackage{hyperref}
\usepackage[T1]{fontenc}
\setlist[itemize]{leftmargin=1.4em, itemsep=0.25em, topsep=0.3em}

\title{Clean Air Act: Statutory Overview}
\author{}
\date{}

\begin{document}
\maketitle

\noindent This overview states only facts confirmed against EPA documentation (sources listed at the
end). It describes the statutory architecture of the Clean Air Act (CAA); for the current status of
specific rules and standards, consult current EPA sources.

\section*{The statute}
The Clean Air Act is codified at 42 U.S.C. \S7401 et seq. EPA's summary notes the Act ``was amended in
1977 and 1990''; the NAAQS program reflects amendments through 1990.

\section*{National Ambient Air Quality Standards (NAAQS)}
The Clean Air Act requires EPA to set National Ambient Air Quality Standards (40 CFR part 50) for six
principal (``criteria'') air pollutants:
\begin{itemize}
  \item carbon monoxide, lead, nitrogen dioxide, ozone, particulate matter, and sulfur dioxide.
\end{itemize}
EPA sets two kinds of standards:
\begin{itemize}
  \item \textbf{Primary} standards ``provide public health protection, including protecting the health of
  `sensitive' populations such as asthmatics, children, and the elderly.''
  \item \textbf{Secondary} standards ``provide public welfare protection, including protection against
  decreased visibility and damage to animals, crops, vegetation, and buildings.''
\end{itemize}
NAAQS concern ambient air quality in a geographic area rather than any individual source.

\section*{State Implementation Plans (SIPs)}
Under Section 110 of the Clean Air Act, each state must adopt and submit to EPA for approval a SIP
providing for the maintenance, implementation, and enforcement of the NAAQS, including a permit program
for the construction and modification of stationary sources. SIP requirements are federally enforceable
under Section 113 of the Act (reference 40 CFR Part 52). Where a state does not produce an adequate plan,
a Federal Implementation Plan (FIP) may be used.

\section*{Hazardous air pollutants (Section 112)}
Section 112 of the Clean Air Act addresses emissions of hazardous air pollutants (HAPs) --- pollutants
associated with serious health effects, including cancer. EPA currently lists 188 hazardous air
pollutants. Section 112 directs EPA to set technology- and performance-based standards --- ``maximum
achievable control technology'' (MACT) standards, issued as National Emission Standards for Hazardous Air
Pollutants (NESHAP). MACT standards are promulgated under 40 CFR Part 63.

\section*{Source and permitting programs}
The following CAA programs are recorded as regulatory air programs in EPA's facility systems:
\begin{itemize}
  \item \textbf{Prevention of Significant Deterioration (PSD).} Part C of the Act sets requirements to
  prevent significant deterioration of air quality in areas designated attainment or unclassifiable, and
  requires a major emitting facility to install best available control technology (per Section 165).
  \item \textbf{Title V operating permits.} Under 40 CFR Part 70, the operating-permit program requires
  all major stationary sources (and certain other sources) to hold a comprehensive operating permit;
  Title V does not impose new requirements but provides a single permit assuring compliance with
  applicable requirements, with a duration not to exceed five years.
  \item \textbf{Acid Rain Program.} Requires major reductions of sulfur dioxide and nitrogen oxide
  emissions from electric utilities through a ``cap and trade'' approach, with continuous emissions
  monitoring of affected sources.
  \item \textbf{New Source Review (NSR) and New Source Performance Standards (NSPS).} Recorded as CAA air
  programs applicable to stationary sources.
  \item \textbf{Federally Enforceable State Operating Permit (FESOP).} Establishes federally enforceable
  limits on a source's potential to emit, keeping it below the Title V applicability threshold (used for
  sources classified synthetic minor).
  \item \textbf{Stratospheric ozone protection (Title VI).} EPA programs protecting the ozone layer,
  including CFC tracking (reference 40 CFR Part 82).
\end{itemize}

\section*{Sources}
\begin{itemize}
  \item EPA, \emph{Summary of the Clean Air Act}, \url{https://www.epa.gov/laws-regulations/summary-clean-air-act}
  \item EPA, \emph{Criteria Air Pollutants}, \url{https://www.epa.gov/criteria-air-pollutants}
  \item EPA, \emph{NAAQS Table}, \url{https://www.epa.gov/criteria-air-pollutants/naaqs-table}
  \item EPA, \emph{What are Hazardous Air Pollutants?}, \url{https://www.epa.gov/haps/what-are-hazardous-air-pollutants}
  \item EPA, \emph{AFS Data Download} documentation (April 2014) --- program descriptions for SIP, PSD,
  NESHAP/MACT, Title V, Acid Rain, FESOP, and CFC tracking.
\end{itemize}

\end{document}
